% !TEX program = xelatex
\documentclass[11pt,a4paper]{article}

\usepackage{fontspec}
\usepackage{polyglossia}
\usepackage{amsmath,amssymb,amsfonts}
\usepackage{xcolor}
\usepackage{graphicx}
\usepackage{geometry}

\geometry{left=2cm,right=2cm,top=2cm,bottom=2cm}

\setdefaultlanguage{persian}
\setotherlanguage{english}
\setmainfont{XB Niloofar}

\newcommand{\en}[1]{\textenglish{#1}}
\newcommand{\separator}{\vspace{1em}\noindent\rule{\linewidth}{0.4pt}\vspace{1em}}
\newcommand{\thickseparator}{\vspace{1em}\noindent\rule{\linewidth}{1pt}\vspace{1em}}

\title{\textbf{پاسخ سوال ۲} \\[0.4cm] \small رگرسیون لجستیک برای تحلیل احساسات}
\author{محمد باغبان زاده \\ شماره دانشجویی: شناسهٔ دانشجو}
\date{}
\linespread{1.35}

\begin{document}

\maketitle
\thickseparator

\section{مفاهیم پایه: رگرسیون لجستیک برای تحلیل احساسات}

\subsection{رگرسیون لجستیک و تابع سیگموئید}
در مسائل رده‌بندی دودویی (مانند تشخیص مثبت یا منفی بودن احساسات یک متن)، مدل رگرسیون لجستیک احتمال تعلق نمونه به کلاس مثبت (معمولاً \(y=1\)) را محاسبه می‌کند. این کار با استفاده از \textbf{تابع سیگموئید} انجام می‌شود:

\[
\sigma(z) = \frac{1}{1 + e^{-z}}
\]

که در آن \(z\) ترکیب خطی ویژگی‌ها و وزن‌ها به اضافهٔ بایاس است:

\[
z = w \cdot x + b = \sum_{j=1}^{n} w_j x_j + b
\]

بنابراین خروجی مدل (احتمال کلاس مثبت) برابر است با:

\[
\hat{y} = P(y=1|x) = \sigma(w \cdot x + b)
\]

\subsection{آستانه تصمیم‌گیری \en{(Decision Boundary)}}
برای تبدیل این احتمال به یک تصمیم رده‌بندی، از یک \textbf{آستانه} استفاده می‌شود. رایج‌ترین مقدار آستانه، \(0.5\) است:

\[
\text{کلاس پیش‌بینی‌شده} = 
\begin{cases} 
1 & \text{\en{if }} \hat{y} \ge 0.5 \\
0 & \text{\en{if }} \hat{y} < 0.5
\end{cases}
\]

این قاعده در بخش \textbf{۴.۳ کتاب} تحت عنوان \en{Classification with Logistic Regression} توضیح داده شده است. در مثال کتاب (شکل ۴.۲)، برای یک نمونه با \(\hat{y}=0.70\)، تصمیم به کلاس \(1\) (احساس مثبت) گرفته شده است.

\subsection{تابع هزینهٔ کراس-آنتروپی \en{(Cross-Entropy Loss)}}
برای اندازه‌گیری خطای مدل، از \textbf{تابع هزینهٔ کراس-آنتروپی} استفاده می‌شود. این تابع در بخش \textbf{۴.۵ کتاب} به تفصیل معرفی شده است. فرمول آن برای یک نمونه \((x, y)\) به صورت زیر است:

\[
L_{CE}(\hat{y}, y) = -[y \log \hat{y} + (1-y) \log(1-\hat{y})]
\]

\begin{itemize}
    \item اگر \(y=1\) (نمونه مثبت)، فرمول ساده می‌شود به \(L = -\log \hat{y}\)
    \item اگر \(y=0\) (نمونه منفی)، فرمول ساده می‌شود به \(L = -\log(1-\hat{y})\)
\end{itemize}

همان‌طور که در مثال کتاب دیده می‌شود، هرچه \(\hat{y}\) به \(y\) نزدیک‌تر باشد، مقدار loss کوچک‌تر خواهد بود. برای نمونهٔ کتاب با \(\hat{y}=0.70\) و \(y=1\)، مقدار loss برابر \(-\log(0.70) \approx 0.36\) محاسبه شده است.

\thickseparator

\section{حل گام‌به‌گام مسأله}

\subsection{داده‌های مسأله}
\begin{align*}
w &= [1.2,\ -0.8,\ 2.5,\ -1.1] \\
b &= 0.3 \\
x &= [2,\ 1,\ 0,\ 3] \quad \text{(شمارش وقوع کلمات: \en{terrible, masterpiece, excellent, boring})}
\end{align*}

\subsection{گام ۱: محاسبهٔ خروجی مدل (احتمال کلاس مثبت)}
ابتدا مقدار \(z\) را محاسبه می‌کنیم:

\[
z = w \cdot x + b = (1.2 \times 2) + (-0.8 \times 1) + (2.5 \times 0) + (-1.1 \times 3) + 0.3
\]

\[
z = 2.4 - 0.8 + 0 - 3.3 + 0.3 = (2.4 - 0.8) = 1.6, \quad 1.6 - 3.3 = -1.7, \quad -1.7 + 0.3 = -1.4
\]

\[
z = -1.4
\]

حال تابع سیگموئید را اعمال می‌کنیم:

\[
\hat{y} = \sigma(-1.4) = \frac{1}{1 + e^{-(-1.4)}} = \frac{1}{1 + e^{1.4}}
\]

مقدار \(e^{1.4}\) را تقریب می‌زنیم (\(e^{1.4} \approx 4.0552\)):

\[
\hat{y} \approx \frac{1}{1 + 4.0552} = \frac{1}{5.0552} \approx 0.1978
\]

\subsection{گام ۲: تعیین کلاس نمونه با آستانه \(0.5\)}
از آنجا که \(\hat{y} \approx 0.1978\) و این مقدار از \(0.5\) کوچک‌تر است، طبق قاعدهٔ تصمیم‌گیری:

\[
\text{کلاس پیش‌بینی‌شده} = 0
\]

بنابراین مدل این سند را به \textbf{کلاس منفی (احساس منفی)} تعلق می‌دهد.

\separator
\subsection{گام ۳: محاسبهٔ مقدار Loss با فرض برچسب واقعی منفی (\(y=0\))}
در سوال گفته شده: «اگر برچسب واقعی این سند منفی است». بنابراین \(y = 0\).

برای \(y=0\)، فرمول کراس-آنتروپی ساده می‌شود به:

\[
L_{CE}(\hat{y}, y=0) = -\log(1-\hat{y})
\]

مقدار \(\hat{y} \approx 0.1978\) را جایگذاری می‌کنیم:

\[
1 - \hat{y} = 1 - 0.1978 = 0.8022
\]

\[
L = -\log(0.8022)
\]

با استفاده از لگاریتم طبیعی (همان‌طور که در مثال کتاب از \(\log\) طبیعی استفاده شده است):

\[
\log(0.8022) \approx -0.2204
\]

بنابراین:

\[
L = -(-0.2204) = 0.2204
\]

\separator
\section{جمع‌بندی و پاسخ نهایی}

\begin{itemize}
    \item \textbf{کلاس پیش‌بینی‌شده برای نمونه \(x = [2,1,0,3]\):}  
    با توجه به \(\hat{y} \approx 0.1978 < 0.5\)، نمونه به \textbf{کلاس ۰ (منفی)} تعلق دارد.
    
    \item \textbf{مقدار Loss با فرض برچسب واقعی \(y = 0\):}  
    \[
    L_{CE} \approx 0.2204
    \]
\end{itemize}

% \thickseparator
% \begin{center}
% \fbox{\parbox{0.9\linewidth}{\centering
% \textbf{پاسخ نهایی:} \\
% ۱) کلاس پیش‌بینی‌شده: \(0\) (منفی) \\
% ۲) مقدار \en{Loss}: \(0.2204\)
% }}
% \end{center}

\thickseparator

\end{document}